\documentclass[12pt]{article}
\usepackage[utf8]{inputenc}
\usepackage[T1]{fontenc}
\usepackage[margin=1in]{geometry}
\usepackage{setspace}
\usepackage{xcolor}
\usepackage[colorlinks=true, urlcolor=blue!60!black, linkcolor=blue!60!black]{hyperref}
\setlength{\parindent}{2em}
\setlength{\parskip}{1.25em}
\renewcommand{\baselinestretch}{1.0}
\begin{document}
\begin{center}
{\LARGE \textit{R.A.V. v. City of St. Paul}, 505 U.S. 377 (1992)\par}
\vspace{0.8em}
{\large Case Brief --- Four-Page Edition\par}
\end{center}
\vspace{1em}

\section{Memory Jogger}

After a teenager was charged with burning a cross in a Black family's yard, the Court unanimously struck down St.~Paul's bias-motivated crime ordinance. Justice Scalia's five-Justice majority held that even within an unprotected category like fighting words, government may not single out disfavored topics or viewpoints for punishment (\href{https://supreme.justia.com/cases/federal/us/505/377/}{505 U.S. 377}).

\section{Detailed Case Facts}

Before dawn on June~21, 1990, R.A.V., then a juvenile, and several other teenagers allegedly taped broken chair legs into a crude cross and burned it inside the fenced yard of a Black family who lived across the street from where R.A.V. was staying in St.~Paul, Minnesota. Although content-neutral charges such as terroristic threats, arson, and criminal damage to property were available, the City charged R.A.V. under its Bias-Motivated Crime Ordinance. The ordinance made it a misdemeanor to place on public or private property ``a symbol, object, appellation, characterization or graffiti, including, but not limited to, a burning cross or Nazi swastika, which one knows or has reasonable grounds to know arouses anger, alarm or resentment in others on the basis of race, color, creed, religion or gender.''

\section{Procedural History}

The trial court dismissed the charge as overbroad and content-based. The Minnesota Supreme Court reversed, construing the ordinance to reach only \textit{Chaplinsky} fighting words and holding that any content distinction was narrowly tailored to a compelling interest in protecting the community against bias-motivated threats. The Supreme Court granted certiorari and reversed on June~22, 1992.

\section{Judicial Votes}

Reversed and remanded, \textbf{9--0} in judgment, \textbf{5--4} on rationale.
\begin{itemize}
\item \textbf{Majority:} Justice Scalia, joined by Chief Justice Rehnquist and Justices Kennedy, Souter, and Thomas.
\item \textbf{Concurring in the judgment:} Justice White, joined by Justices Blackmun and O'Connor, and by Justice Stevens except Part~I-A; Justice Blackmun; Justice Stevens, joined in Part~I by Justices White and Blackmun.
\item \textbf{Dissenting:} None.
\end{itemize}

\section{Holding}

Even as narrowed to proscribable fighting words, the ordinance is facially unconstitutional because it discriminates by content and viewpoint within a proscribable category and is not necessary to serve a compelling interest. Content distinctions within such a category are permissible only when they track the very reason the category is proscribable, target secondary effects or conduct, or present no realistic possibility of official suppression of ideas.

\section{Analysis of Opinions}

\textbf{Majority (Scalia).} Accepting the state court's construction, Scalia bypassed overbreadth. Obscenity, defamation, and fighting words may be regulated for their proscribable content, but they are not invisible to the Constitution. Like a noisy sound truck, fighting words are an unprotected mode of expression, yet government may not regulate them out of hostility to their message, as by banning only libel critical of government. Scalia identified three exceptions: distinctions based on the very reason the class is proscribable (such as threats against the President); secondary effects under \textit{Renton}, or general conduct rules like Title~VII; and distinctions posing no realistic possibility that suppression of ideas is afoot.

St.~Paul's ordinance reached fighting words on five disfavored topics while ignoring equally vicious ones about political affiliation, union membership, or sexual orientation. In practice it discriminated by viewpoint: advocates of tolerance could use other fighting words freely, while their opponents could not. None of the exceptions applied, because the City targeted the message's emotive impact. Protecting historically discriminated-against groups is compelling, but a content-neutral ban on all fighting words would serve that interest equally well.

\textbf{White (concurring in the judgment).} The ordinance was simply overbroad, since it still reached speech causing only anger or resentment. White attacked the majority's new underbreadth theory as abandoning the categorical approach and threatening Title~VII.

\textbf{Blackmun (concurring in the judgment).} The majority either weakened scrutiny generally or issued an aberration driven by opposition to ``politically correct speech.''

\textbf{Stevens (concurring in the judgment).} Rejecting both rigid approaches, Stevens favored a contextual inquiry and would have upheld a non-overbroad ordinance barring ``punches below the belt.''

\textbf{Later developments.} \textit{Wisconsin v.\ Mitchell} (1993) upheld bias-crime sentence enhancements as regulating conduct, and \textit{Virginia v.\ Black} (2003) upheld bans on cross burning with intent to intimidate under the first exception.

\section{Examples: Future Applications}

\begin{enumerate}
\item \textbf{Same-Side (Invalid): Insulting the Police.} A city bans only fighting words that disparage law enforcement. It discriminates by viewpoint on its face and fails strict scrutiny, since a neutral ban would protect officers.
\item \textbf{Same-Side (Invalid): Obscenity Mocking Officials.} An enhancement for obscene videos that ridicule elected officials targets political criticism, not what makes obscenity unprotected.
\item \textbf{Opposite-Side (Valid): Threats Against Election Workers.} Heavier penalties for true threats against poll workers fit the first exception, like threats against the President, and favor no viewpoint.
\item \textbf{Opposite-Side (Valid): Bias Enhancement for Conduct.} Charging R.A.V. with content-neutral crimes plus a victim-selection enhancement raises no problem under \textit{Wisconsin v.\ Mitchell}.
\item \textbf{Fence-Sitter: A Campus Harassment Act.} A ban on campus fighting words that demean on the basis of nationality, religion, or political ideology is topic-based like St.~Paul's, but the State may claim the first exception and evenhandedness across viewpoints.
\end{enumerate}

\section{Critique}

\textbf{The paradox of underinclusion.} Critics, echoing White and Stevens, note a striking paradox: a city that bans all fighting words suppresses more speech yet is safer than one targeting only racial slurs, turning the preference for narrow tailoring upside down. The line between the first exception and forbidden content discrimination is also elastic, as \textit{Virginia v.\ Black}'s approval of a targeted cross-burning ban shows.

\textbf{The defense.} Defenders answer that \textit{R.A.V.} keeps majorities from weaponizing broad low-value categories to silence disfavored political or social viewpoints.

\section{Key Quotations}

\begin{quote}
``What they mean is that these areas of speech can, consistently with the First Amendment, be regulated because of their constitutionally proscribable content (obscenity, defamation, etc.)---not that they are categories of speech entirely invisible to the Constitution, so that they may be made the vehicles for content discrimination unrelated to their distinctively proscribable content.''
\end{quote}
\noindent --- Justice Scalia, Opinion of the Court, 505~U.S.~at~383--384.

\begin{quote}
``St. Paul has no such authority to license one side of a debate to fight freestyle, while requiring the other to follow Marquis of Queensberry rules.''
\end{quote}
\noindent --- Justice Scalia, Opinion of the Court, 505~U.S.~at~392.

\begin{quote}
``Let there be no mistake about our belief that burning a cross in someone's front yard is reprehensible. But St. Paul has sufficient means at its disposal to prevent such behavior without adding the First Amendment to the fire.''
\end{quote}
\noindent --- Justice Scalia, Opinion of the Court, 505~U.S.~at~396.

\end{document}
